\documentclass[10pt]{article}
\usepackage[margin=20mm]{geometry}
\usepackage{amsmath,amssymb,mathtools,fontspec,unicode-math,xurl,hyperref}
\IfFontExistsTF{DejaVu Serif}{\setmainfont{DejaVu Serif}}{\setmainfont{Latin Modern Roman}}
\setmathfont{latinmodern-math.otf}
\setlength{\parindent}{0pt}\setlength{\parskip}{6pt}\sloppy
\title{A version-specific review of a corrected
potential-complexity criterion}
\author{AI-assisted technical review}
\date{}
\begin{document}
\maketitle
Abstract: The publisher version of Martino Lupini’s paper contains a false equivalence between potential \(\Sigma^0_2\) complexity of a coset relation and \(D(\Pi^0_2)\) membership of its subgroup. We verify a complete counterexample from the author’s dated corrigendum, including Polishability, the exact Borel upper bound, failure of the \(F_\sigma\) property, and the potential-complexity obstruction. The printed proof uses a hyperspace-valued map without establishing that its values are closed in the refined topology. This is a Major defect in the bound original version. The correction is already author-documented; it is not claimed as a new discovery. The wider replacement classification is recorded with a separate verification limit. Scientific adjudication is pending.\par
\section*{Version and scope}
The primary source is M. Lupini, Complexity classes of Polishable subgroups, J. Eur. Math. Soc. 28 (2026), 2197–2232, DOI 10.4171/JEMS/1605 (https://doi.org/10.4171/JEMS/1605). The acquired 36-page publisher PDF has SHA-256 beginning `cafc0b5fb2d8c752`; PDF page \(j\) is printed page \(2196+j\). Its first page carries copyright 2025. The frozen project identifier M25-0232 and year mapping are retained.\par
The separately acquired Corrigendum to Complexity Classes of Polishable Subgroups, dated 10 September 2026, is Zenodo version 3, DOI 10.5281/zenodo.22686173 (https://doi.org/10.5281/zenodo.22686173). Its 35-page PDF has SHA-256 beginning `c64a640eef9ba856`. The public record confirms that date and version. The original and corrigendum are different scientific objects; no replacement of the frozen source is implied.\par
Close reading covers the introductory statements and definitions, original Theorem 3.3 and its proof, the first Solecki-stage description, and Lemmas 6.2–6.8 with their immediate dependencies. Corrigendum pp. 1–12 were read, concentrating on Proposition 2, Example 3, the diagnosis on pp. 8–9 and the restricted replacement on pp. 9–11. The later transfinite operator-range constructions were not independently re-proved. This scope does not certify all original or corrected classifications.\par
\section*{The false assertion and its hypotheses}
Here \(D(\Pi^0_2)\) means differences of two \(G_\delta\) sets, equivalently an intersection of a \(G_\delta\) set with an \(F_\sigma\) set. A Polishable subgroup has its unique compatible Polish group topology, which may be finer than the ambient subspace topology. Potential complexity refers to membership after a Borel-preserving Polish refinement, not an assertion of an exact minimal class.\par
Problem 1.  Original Theorem 3.3(2), p. 2201 (PDF p. 5), asserts that \(E_H^G\) is potentially \(\Sigma^0_2\) if and only if \(H\in D(\Pi^0_2)(G)\). Its converse is false. The potential-\(\Sigma^0_2\) assertion of Lemma 6.5, p. 2208 (PDF p. 12), is also false.\par
We reproduce and verify the author’s Example 3 rather than infer an error solely from the existence of a corrigendum. Put \(G=\mathbb R^{\mathbb N}\) with its product topology and \[H=bv_0(\mathbb R)=
 \left\{x:x_n\longrightarrow0,\quad
              \sum_{n\ge1}|x_n-x_{n+1}|<\infty\right\}.\] This is an additive subgroup. With norm \(\|x\|_{\rm var}=\sum_{n\ge1}|x_n-x_{n+1}|\), it is isometrically isomorphic to \(\ell^1\) via \[(a_n)\longmapsto\left(\sum_{k\ge n}a_k\right)_n,
 \qquad a_n=x_n-x_{n+1}.\] The inverse formula follows by telescoping and the limit-zero condition. Also \(|x_n|\leq\|x\|_{\rm var}\), so inclusion into \(G\) is continuous. Thus \(H\) is a Polishable subgroup of the Polish group \(G\).\par
Let \(BV\) be the set of sequences of finite variation. For every integer \(r\ge1\), \[F_r=\bigcap_{N\ge1}
 \left\{x:\sum_{n=1}^{N}|x_n-x_{n+1}|\leq r\right\}\] is product-closed, so \(BV=\bigcup_r F_r\) is \(F_\sigma\). Every finite-variation sequence converges. The set \[L=\{x:\liminf_n|x_n|=0\}
   =\bigcap_{k,m\ge1}\ \bigcup_{n\ge m}\{x:|x_n|<1/k\}\] is \(G_\delta\). Consequently \(H=BV\cap L\in D(\Pi^0_2)(G)\); in particular \(H\in\Sigma^0_3(G)\) as required by Lemma 6.5.\par
\subsection*{Why this subgroup is not \(F_\sigma\)}
Suppose \(H=\bigcup_j C_j\) with \(C_j\) product-closed in \(G\). Their inverse images in the Banach topology of \(H\) are closed. Baire category gives \(z\in H\) and \(\varepsilon>0\) such that \(z+B_{\rm var}(0,\varepsilon)\subseteq C_j\) for some \(j\). Then \(C_j-z\) is product-closed, is contained in \(H\), and contains this norm ball. Choose \(0<c<\varepsilon\) and define \[x^{(m)}_n=
 \begin{cases}c,&n\leq m,\\0,&n>m.\end{cases}\] Each \(x^{(m)}\) has variation norm \(c\) and belongs to that ball. Coordinatewise it converges to the nonzero constant sequence \((c,c,\ldots)\), which is not in \(H\). This contradicts product-closedness. Thus \(H\notin\Sigma^0_2(G)\).\par
\subsection*{Why the coset relation is not potentially \(\Sigma^0_2\)}
The necessary criterion can be proved without assuming the disputed theorem. For a Borel subgroup \(K\) of a Polish group \(P\), suppose \(E_K^P\) Borel-reduces to an \(F_\sigma\) equivalence relation via \(f:P\to Y\). There is a comeager \(G_\delta\) set \(C\) on which \(f\) is continuous. Hence for closed sets \(D_j\subseteq P^2\), \[E_K^P\cap C^2=\bigcup_j(D_j\cap C^2).\] Use the convention \(xE_K^Py\iff x^{-1}y\in K\); inversion interchanges this with the original paper’s convention. For \(h\in P\) put \(A_j(h)=\{x:(x,xh)\in D_j\}\). These sets are closed. On the comeager set \(C\cap Ch^{-1}\), their union contains every point if \(h\in K\) and no point otherwise. If \((V_i)\) is a countable basis of nonempty open sets of \(P\), Baire category gives \[K=\bigcup_{j,i}\{h:\ \forall x\in V_i,\ (x,xh)\in D_j\}.\] Each set in this union is closed, as an intersection of closed inverse images. Therefore potential \(\Sigma^0_2\) implies \(K\in\Sigma^0_2(P)\). Conversely, the map \((x,y)\mapsto x^{-1}y\) immediately gives the reverse implication. This verifies the corrected criterion in the author’s Proposition 2.\par
Applied to \(H=bv_0(\mathbb R)\), the necessary direction and the preceding non-\(F_\sigma\) proof show that \(E_H^G\) is not potentially \(\Sigma^0_2\). All hypotheses of both disputed assertions are satisfied. No finite simulation or conjectural complexity lower bound is being substituted for proof.\par
\section*{Proof diagnosis, impact, and repair boundary}
In the original proof of Lemma 6.5, p. 2209 (PDF p. 13), a refined Polish topology \(t\) makes the left \(H\)-action continuous and certain sets \(U_n\) closed. The map \(x\mapsto(U_nx)_n\) is then treated as taking values in the hyperspace of \(t\)-closed sets. The stated properties do not establish that the right translates \(U_nx\) are \(t\)-closed for arbitrary \(x\in G\). The subsequent calculation of inverse images of hit sets addresses measurability and does not supply this missing closed-valuedness. The author’s diagnosis on corrigendum pp. 8–9 agrees with this reading. Requiring joint continuity of the right action of the original group would force the refinement to be the original topology by the orbit map at the identity; that condition cannot simply be inserted.\par
The supported severity is Major: an asserted equivalence between two complexity invariants is false, and its central proof lemma is invalid. Theorem 3.3 explicitly invokes Lemma 6.5. Lemmas 6.7 and 6.8 on p. 2210 also invoke it in the classification argument. This dependency does not mean that every dependent conclusion is false.\par
For this criterion, the verified repair is to replace \(D(\Pi^0_2)\) by \(\Sigma^0_2\) in Theorem 3.3(2). The example does not refute the separate subgroup conclusion \(H\in D(\Pi^0_2)\) of Lemma 6.5: our example satisfies it. The author refutes that separate unrestricted conclusion using a different operator construction, and gives the valid \(\Delta^0_3\) restriction. The restricted proof and the distinction between Lemmas 6.7 and 6.8 were read; the later operator constructions were not completely audited here.\par
The corrigendum also reports changes to Theorems 1.1, 1.3, 1.4, 6.1 and 10.2, while retaining the non-Archimedean list of Theorem 1.2. These are recorded as author-documented broader corrections, not additional independently established findings of this report. No blanket certification of the 35-page correction is intended.\par
\section*{Conclusion}
Dimension | Recorded state\par
Confirmed issue | False potential-\(\Sigma^0_2\) criterion in the exact publisher version\par
Impact | Major; theorem-level assertion and core proof lemma\par
Decisive evidence | Polishable \(bv_0(\mathbb R)\in D(\Pi^0_2)\setminus\Sigma^0_2\) and a proved necessary criterion\par
Repair status | Criterion corrected and verified above; wider classification not fully re-proved\par
Author status | Existing dated corrigendum v3; no new author communication\par
Scientific decision | Pending; human reviews empty\par
\section*{Supporting checks}
The proof above is analytic; no author code or artificial finite test was run. Exact source, TeX, compiler, dependency, PDF and actual page-view bindings are retained under `verification/0232`. The acquired arXiv manuscript remains a separate source version. The bounded supplement search located the author’s corrigendum and its TeX; no dedicated computational package was identified. Absence of code is not a mathematical defect.\par
All new review reasoning and actual compilation/inspection were performed by this assigned web executor. The visible UI/model observation is attributed to Main; backend equivalence was not independently verified. This is an initial review, not independent acceptance.\par
9 M. Lupini, Complexity classes of Polishable subgroups, J. Eur. Math. Soc. 28 (2026), 2197–2232. doi:10.4171/JEMS/1605. Exact acquired publisher PDF. M. Lupini, Corrigendum to Complexity Classes of Polishable Subgroups, version 3, 10 September 2026. doi:10.5281/zenodo.22686173. Exact acquired 35-page PDF; especially Proposition 2, Example 3 and the proof diagnosis on pp. 8–9.\par
\end{document}
